\scanpage{421}
\mspage{213}
% Main text, approximate manuscript lines 1--4.
\deleted{Let} Is uniqueness at \unsure{all|all} events clear?
\deleted{\unsure{Choosing a}} Let us consider the extension \inserted{$\mathcal{E}$} of
$\Gamma$ by $L=\mathrm{Ker}((N_a)_{I'}\to\unsure{N_a|N_0})$,
inverse \unsure{image} of the extension $(N_a)_{I'}$
(of $\unsure{N_{a'}|N_{0'}}$ by $L$) via $\Gamma\to\unsure{N_{a'}|N_{0'}}$, we
% Main text, approximate manuscript lines 5--8.
have a \inserted{partial} splitting of this extension \unsure{over}
the subgroup $\Gamma'$ of $\Gamma$, and the assertion is
that this splitting \uline{\unsure{extends}}, \uline{in a way} \uline{that is unique}, to $\Gamma$.
We may suppose $\Gamma'$ normal in $\Gamma$, and we
% Main text, approximate manuscript lines 9--13.
identify $\Gamma'$ with a subgroup of $\mathcal{E}$. \unsure{All}
\deleted{\unsure{sections}} the sections of $\mathcal{E}$ over $\Gamma$ \unsure{are identified}
with \inserted{\unsure{section}} subgroups $\widetilde{\Gamma}$ of $\mathcal{E}$.
For such a $\widetilde{\Gamma}$, we have of course
$\widetilde{\Gamma}\subset\mathrm{Norm}_{\mathcal{E}}(\Gamma')$
moreover \deleted{\unsure{if we} \illegible{2}} \inserted{we obviously have}
% Main text, approximate manuscript line 14.
\[
\mathrm{Norm}_{\mathcal{E}}(\Gamma')\cap L=L^{\Gamma'} \tag{31}
\]
% Main text, approximate manuscript lines 15--21.
and I presume we must have
$L^{\Gamma'}=1$ (which generalizes condition
\textcircled{3} above\ldots) Now \uline{this condition}
\uline{implies uniqueness} -- namely
$\widetilde{\Gamma}=\mathrm{Norm}_{\mathcal{E}}(\Gamma')$ --
and existence \uline{means} that
$\mathrm{Norm}_{\mathcal{E}}(\Gamma')\to\Gamma$
is an epimorphism (hence an isomorphism\ldots).

% Main text, approximate manuscript lines 22--23.
But we will have to try to make precise,
in certain contexts (at least that of
\scanpage{422}
% Main text, approximate manuscript lines 1--4; continuation of scan 421.
arithmetically faithful actions, i.e. $\Gamma\to\boldsymbol{\Gamma}_a$
injective -- which \unsure{corresponds} \unsure{normally|presumably}
to the case of algebraic curves \unsure{defined} over
algebraic extensions of $\mathbb{Q}$ -- the notion
% Main text, approximate manuscript lines 5--9.
of an ``admissible'' action. For this, we
must return to the relation between
(potential) curves and finite coverings,
which also gives a way of
varying $g$.

% Main text, approximate manuscript lines 10--17.
But I feel like first taking up again
\unsure{under|on} another aspect \inserted{(perhaps more general)} the preceding
formalism, which \unsure{can} perhaps also \unsure{be} applied to the case of \uline{discrete} groups
with loops (I am thinking of the formalism
of the $\unsure{U^!|\mathcal{U}^!}$, linked to outer actions of
finite groups on such \unsure{profinite groups} with loops).

% Main text, approximate manuscript lines 18--20.
Let $\mathcal{X}$ be a set \inserted{$\ne\varnothing$}
(\inserted{\unsure{morally}, a} set of the
``points'' of an \deleted{surface} algebraic curve,
or of a topological surface\ldots), \deleted{and $\Gamma$}
\scanpage{423}
\mspage{214}
% Main text, approximate manuscript lines 1--7.
\deleted{a group operating on $\mathcal{X}$.} We suppose
given, for every finite subset $I$ of $\mathcal{X}$,
with complement $U_I=\mathcal{X}\setminus I$, a
\unsure{groupoid} $\Pi_{U_I}$, having $U_I$ as its set
of \unsure{objects}. Moreover, for $J\supset I$ i.e.
$U_J\subset U_I$, we suppose given a homomorphism
of groupoids
% Main text, approximate manuscript line 8.
\[
\Pi_{U_J}\to\Pi_{U_I}
\]
% Main text, approximate manuscript lines 9--14.
which on objects is the inclusion $U_J\unsure{\subset|\hookrightarrow}\unsure{U_J|U_I}$.
We suppose transitivity (strict) i.e. \deleted{\unsure{which}}
a covariant functor \inserted{from the category} of subsets $U$ of
$\mathcal{X}$ complementary to finite subsets \inserted{(with the inclusions)}, to
\unsure{the} category of groupoids, which on \unsure{the set}
of objects coincides with the obvious functor\ldots

% Left marginal note, running beside approximate lines 10--22.
\margin{If $\mathcal{X}$ is finite and $I=\mathcal{X}$, we
\unsure{suppose} \inserted{\uuline{\unsure{all the same|however}}}
$\Pi_{U_I}\;(=\Pi_{\varnothing})$ \unsure{is}
not the empty groupoid, but a groupoid
\deleted{of any kind} (which \unsure{maps} into the preceding ones\ldots)
but we will not be able to require
\unsure{the} strict transitivity.
We \unsure{will need|would need} \inserted{\unsure{perhaps}} \illegible{1} anabelian \unsure{hypotheses}
on $\Pi_{U_I}$ in this case\ldots}

% Main text, approximate manuscript lines 15--18.
We suppose moreover that for every \deleted{$I\in\mathcal{P}_f(\mathcal{X})$}
$i\in\mathcal{X}$, we give ourselves a \deleted{\unsure{groupoid} \illegible{1}}
$\Pi_i$ (or $\Pi_{D_i^*}$),
and \deleted{\unsure{for every}} for $i\in I\in\mathcal{P}_f(\mathcal{X})$,
a homomorphism of groupoids
% Main text, approximate manuscript line 19.
\[
\Pi_{D_i^*}\to\Pi_{U_I}
\]
% Main text, approximate manuscript lines 20--23.
compatible with the transition morphisms
\inserted{$\Pi_{U_I}\to\Pi_{U_J}$ $(I\supset J)$}.
We suppose that for fixed $I$, we find
on $\Pi_{U_I}$ a structure of groupoid
with loops by
\[
\Pi_{D_I^*}=\coprod_{i\in I}\Pi_{D_i^*}\to\Pi_{U_I}.
\]
\scanpage{424}
% Main-text lines 1--2; insertion above line 2.
without moreover excluding the case where $I=\varnothing$, and where \inserted{($\Pi_{\mathcal{X}}$ \unsure{being} connected, say)} \unsure{its|the} genus \uline{is $0$} --- \deleted{\unsure{thus}}
\margin{\unsure{We suppose} \illegible{4} \unsure{that} \illegible{1} $\mathcal{X}$ \illegible{3} \deleted{\unsure{for} $\mathcal{X}$ \illegible{2}} \deleted{\unsure{two points}}}

% Main-text lines 3--8.
We suppose moreover that if $J\supset I$ i.e. $U_J\subset U_I$, \inserted{and for $i\in J\setminus I$,} the composite homomorphism
\[
\Pi_{D_i^*}\longrightarrow\Pi_{U_J}\longrightarrow\Pi_{U_I}
\]
is ``constant'', in such a way that $\Pi_{U_I}$ is \unsure{obtained} from $\Pi_{U_J}$ by ``plugging the holes'' at $J\setminus I$.

% Main-text lines 9--21; the middle of the example is heavily revised.
Up to now, \unsure{all|all} this could be visualized for example by starting from an orientable topological surface \deleted{of the form \unsure{$X\setminus S$}; \illegible{4} of type} $X$, and a subset \inserted{\illegible{1}} $\mathcal{X}$ of $X$ \deleted{\illegible{1}}\inserted{\unsure{meeting}} \deleted{\unsure{every connected component}}\inserted{\unsure{every} connected component}; by taking for $\Pi_I$ the restriction to $\mathcal{X}\setminus I$ of the fundamental groupoid of $X\setminus I$ (if $I\ne\mathcal{X}$; if $\mathcal{X}=I$, which requires $\mathcal{X}$ finite, we will take the fundamental groupoid of $X\setminus I=X\setminus\{\mathcal{X}\}$; which we will have trouble mapping into the others with \uline{strict} transitivity, no matter!

% Main-text lines 22--23; sentence continues on the next scan.
But, we are mainly interested in the case $\mathcal{X}$ infinite. Where we \unsure{take} $X$ an algebraic curve over
\scanpage{425}
\mspage{215}
% Main-text lines 1--3; continuation from scan 424.
an algebraically closed field, $\mathcal{X}$ a subset of $X$ meeting \unsure{every} connected component, and we define the $\Pi_{U_I}$ as before.

% Main-text lines 4--9, with an insertion between lines 6 and 7.
Let now $\Gamma$ be a group which operates on the situation (\uline{strictly}, be it understood) thus on \unsure{the set} $\mathcal{X}$, \deleted{and} the groupoids \unsure{$\Pi_{U_I}$}, \inserted{the $\Pi_{D_i^*}$, (N.B. which it \unsure{permutes} among themselves),} commuting with the homomorphisms \unsure{d|of}
\[
\Pi_{U_I}\longrightarrow\Pi_{U_J}
\]
and the $\Pi_{D_i^*}\longrightarrow\Pi_{U_I}$.
\margin{\deleted{(which it \unsure{permutes} \unsure{among themselves})}}

% Main-text lines 10--13; the wording after Nous is rewritten interlinearly.
(We will have difficulties, if $\mathcal{X}$ is finite, for $\Pi_{U_I}$ when $I=\mathcal{X}$, for the \uline{strict} commutation \deleted{\unsure{of the}} \unsure{$\Pi_{U_{\mathcal{X}}}$}$\longrightarrow\Pi_{U_J}$ \unsure{..} but let us pass on...)
We \deleted{make \unsure{the hypothesis|the hypotheses},} \inserted{suppose defined a notion of \uline{admissible} subgroup of $\Gamma$, (they are in any case $\ne1$) such that \deleted{\illegible{1}} if $\Gamma'$ is admissible \unsure{then} every subgroup $\Gamma''\supset\Gamma'$ also.} \deleted{\unsure{following}} \inserted{We \deleted{\unsure{suppose}}}

% Canceled condition and canceled bracketed development in the lower half.
\deleted{\textcircled{\(1^\circ\)} For $P\in\mathcal{X}$, its stabilizer $\Gamma_P$ in $\Gamma$ is \deleted{$\ne1$}\inserted{admissible} (hence $\ne1$)}
\inserted{\deleted{\unsure{$U_I$ stable under $\Gamma_P$ i.e. $I$ stable}}}

\deleted{\unsure{\textcircled{\(2^\circ\)}} \illegible{1} (Let us note that if $U_I\ni P$, then $P$ is an object of $\Pi_{U_I}$ with an operation of $\Gamma_P$ on it, \deleted{\illegible{1}}\inserted{We} \unsure{find} in this way an \deleted{map}\inserted{operation} of $\Gamma_P$ on \unsure{$\pi_1(U_I,P)$|$\Pi_1(U_I,P)$}, and we suppose}

\deleted{\textcircled{\(2^\circ\)} \unsure{$\pi_1(U_I,P)$|$\Pi_1(U_I,P)$} \deleted{\unsure{call}}\inserted{\unsure{say that a}} subgroup $\Gamma'$ of $\Gamma$ is \textcircled{\(1^\circ\)} If $\Gamma'$ \uline{admissible} if \deleted{\illegible{1}} $\mathcal{X}^{\Gamma'}$ \unsure{meets} \unsure{every} connected component of $\mathcal{X}=U_{\varnothing}$ (which}